\section{From values to points}\label{sec:points}

Convex optimization algorithms are value algorithms at heart. The ellipsoid
method and its relatives return a feasible point whose objective value is
close to the optimum; they say little about where the optimizers are. Many
uses in mixed-integer nonlinear optimization need more. Branching and
variable fixing act on coordinates of an optimizer. A decomposition method
that passes the optimizer of a subproblem to another stage needs the same
optimizer at every requested accuracy; otherwise successive refinements may
move between different optimal points. Exact certification needs labels,
such as the set of active constraints, that no finite accuracy settles.
This section asks when value accuracy can be converted into point accuracy
in the ordinary bit model, and it identifies where the conversion fails.

\paragraph{Output contracts.}
We minimize $f$ over a rational polyhedron $P$, with optimal value $f_*$
and optimal set $S$, and use the output contracts of
Section~\ref{sec:models-outputs}. The requested precision is $q$ bits, and
$\varepsilon=2^{-q}$; running times are polynomial in $q$, not in
$1/\varepsilon$. A \emph{value enclosure} (Definition~\ref{def:models-value})
is a feasible rational $x$ with a rational $\ell$ such that
$\ell\le f_*\le f(x)\le\ell+\varepsilon$. A \emph{set approximation}
(Definition~\ref{def:models-points}) is a feasible rational $x$ with
$\dist(x,S)\le\varepsilon$; approximations returned for different $q$ may
approach different optimizers. A \emph{fixed selector} approximates one
optimizer $p$, determined by the instance alone, at every precision; our
selector is the optimizer of minimum Euclidean norm in the original
coordinates. Beyond these numerical outputs are exact predicates, such as
whether a coordinate of the optimizer equals its bound
(Definition~\ref{def:models-predicates}), exact representations, such as the
minimal polynomial of an optimizer coordinate
(Definition~\ref{def:models-representations}), and certificates, such as a
rational box that encloses the optimizer
(Proposition~\ref{prop:points-rectangle}).

A fixed selector gives set approximations. A set approximation $x$ gives the
value enclosure $[f(x)-\Lambda\varepsilon,f(x)]$ whenever $\Lambda$ bounds
$\norm{\nabla f}$ on a convex region containing $x$ and its nearest
optimizer. Exact predicates and representations are formats that no amount
of numerical precision supplies by itself; exact comparisons are the
subject of Section~\ref{sec:points-limits} and of
Theorem~\ref{thm:exact-upper}. Polynomials are explicit: each nonzero
coefficient is listed with its exponent vector. Convexity is a promise in
the sense of Section~\ref{sec:models-promises}.

\paragraph{Results.}
Lemma~\ref{lem:selector} reduces the fixed selector to an \emph{effective
error modulus}: a bound $\dist(x,S)\le\Gamma(f(x)-f_*)^{1/d}$ whose exponent
and constant are known and whose constant $\Gamma$ has polynomially many
bits. With such a modulus and a bound $R$ on the norm of the selected
optimizer, a value algorithm becomes a fixed-selector algorithm at the price
of $O(d(1+q+\log\Gamma+\log R))$ value bits. The section determines when an
effective modulus exists.
\begin{itemize}
\item For an explicit sparse polynomial convex on all of $\R^n$ and an
arbitrary rational polyhedron, Theorem~\ref{thm:global-point} decides
emptiness and unboundedness, proves attainment, computes a radius for the
minimum-norm optimizer and an effective modulus with exponent $1/D$, and
approximates that optimizer to $2^{-q}$ in time polynomial in $L$, the
numerical degree $D$, and $q$.
\item A polynomial of degree three that is convex only on a bounded polytope
also has an effective modulus, with exponent $1/4$
(Theorem~\ref{thm:cubic-point}).
\item For quartics convex only on a box the conversion fails.
Unconditionally, for an explicit family, converting values into points
through any error modulus needs a number of value bits exponential in the
dimension, and isotropic regularization needs a parameter whose bit length is
exponential in the dimension (Proposition~\ref{prop:points-regularization}).
Conditionally, one constant-accuracy point decides Square Root Sum or
$\PosSLP$ (Theorem~\ref{thm:points-quartic-lower}).
\item Even where points are easy, exact labels
(Proposition~\ref{prop:points-active}), the minimum-norm selector among
several optimizers (Proposition~\ref{prop:points-selector}), expanded
algebraic output (Proposition~\ref{prop:points-algebraic-output}), and
rational box certificates (Proposition~\ref{prop:points-rectangle}) can be
substantially harder than a point. Table~\ref{tab:points-summary} summarizes
the picture.
\end{itemize}

\paragraph{Relation to prior work.}
Weak convex optimization in the bit model is classical
\citep{GroetschelLovaszSchrijver1988}, and Lemma~\ref{lem:convex-value}
applies it with an exact feasibility repair. For polynomials convex on
$\R^n$ and rational polyhedra, \citet[Theorem~1.1 and
Corollary~1.2]{SlotSteurerWiedmer2025} proved that the infimum is either
$-\infty$ or attained at a point whose norm has a polynomially bounded
logarithm, and that unboundedness can be detected and the optimal value
approximated in polynomial time. Their structure theory separates the directions in
which $f$ is invariant from those in which it changes linearly, reading
them from the coefficients of the gradient, and it bounds the remaining
part of $f$ below by a strongly convex quadratic; the underlying tangent
quadratic bound appears in the proof of Lemma~5 of
\citet{AhmadiChaudhryZhang2024}. We use the same structural facts. We
rederive attainment and the radius because the point theorem needs explicit
rational constants for sparse inputs of variable degree, and because
transformed polynomials must be evaluated rather than expanded
(Remark~\ref{rem:points-expansion}). H\"older error bounds for convex
polynomials are classical: \citet[Theorem~4.2 and Corollary~4.1]{Li2010}
and \citet[Theorem~1]{Li2013} prove global bounds with exponent
$1/((D-1)^n+1)$ and a constant whose existence is shown;
\citet{Yang2009} also studies error bounds for convex polynomials, and we
make no claim about the exponents or constants obtained there. We do not
claim the exponent $1/D$ as new. The contribution of
Theorem~\ref{thm:global-point} is a constant computed from the input with
polynomially many bits, valid on the whole, possibly unbounded,
polyhedron, together with its use to approximate one fixed optimizer at
every precision. For degree two, \citet{KozlovTarasovKhachiyan1980} compute
exact rational optimizers, a stronger output. The distinction between
approximating the residual of a solution and approximating the solution
itself, and the hardness of the latter for equilibria, is due to
\citet{EY2010}; Theorem~\ref{thm:points-quartic-lower} exhibits the same
phenomenon for minimizers of convex polynomials.

\subsection{Three tools}\label{sec:points-tools}

The first tool turns explicit convex objectives into value enclosures. It
assumes only convexity on the feasible set and exact evaluation of values
and subgradients.

\begin{lemma}[Convex value interface]\label{lem:convex-value}
Let $P=\{x\in\R^n:Ax\le b\}$ with rational $A,b$, and let $\varrho\ge1$ be a
rational number with $P\subseteq[-\varrho,\varrho]^n$.
\begin{enumerate}[label=(\alph*)]
\item Let $\varphi:P\to\R$ be continuous and convex. Suppose that rationals
$W\ge\max\{1,\sup_P|\varphi|\}$ and $G\ge1$ are given, together with an
algorithm that, for every rational $x\in P$, returns the rational value
$\varphi(x)$ and a rational vector $s(x)$ with $\norm{s(x)}\le G$ and
$\varphi(x')\ge\varphi(x)+s(x)^{\mathsf T}(x'-x)$ for all $x'\in P$, in time
bounded by a fixed polynomial in $\Delta+\langle x\rangle$; here $\Delta$ is
the encoding length of $A,b,\varrho,W,G$ and of the description of
$\varphi$. Then there is an algorithm that, given a rational $\eta>0$,
either reports $P=\emptyset$ or returns a rational $\hat x\in P$ and a
rational $\ell$ with
\[
 \ell\le\min_P\varphi\le\varphi(\hat x)\le\ell+\eta .
\]
Its running time is bounded by a polynomial in $\Delta+\langle\eta\rangle$
that depends only on the polynomial bounding the evaluation time.
\item In particular, let $\varphi(x)=f(Tx+t)+\tau\norm x^2$, where $f$ is an
explicit sparse polynomial of numerical degree at most $D$, the matrix $T$
and vector $t$ are rational, $\tau\ge0$ is rational, and $\varphi$ is convex
on $P$. Then the conclusion of (a) holds in time polynomial in $D$ and the
encoding lengths of $A,b,\varrho,f,T,t,\tau,\eta$. The composition
$f(Tx+t)$ is evaluated at query points and never expanded.
\end{enumerate}
\end{lemma}

The proof in Appendix~\ref{app:points-tools} reduces $P$ to its affine hull
by linear programming, places a rational ball inside the reduced polytope,
applies the weak optimization theorem of
\citet[Corollary~(4.2.7)]{GroetschelLovaszSchrijver1988} to a capped
epigraph, and repairs the approximately feasible output by an explicit
homothety toward the center of the ball. The returned point is exactly
feasible, also when $P$ is lower dimensional. For part (b),
Lemma~\ref{lem:points-sparse} supplies exact values and gradients, and the
bounds $W$ and $G$, by evaluating the sparse polynomial; this matters when
the degree grows, because expanded compositions can have superpolynomially
many monomials (Remark~\ref{rem:points-expansion}). For polynomials convex
on $\R^n$ and unbounded polyhedra, value approximation is due to
\citet{SlotSteurerWiedmer2025}; Theorem~\ref{thm:global-point}(c) recovers
it through Lemma~\ref{lem:convex-value} after a radius computation.

The second tool is a Hoffman bound with an explicit constant
\citep{Hoffman1952}. Its right-hand sides may be irrational. This matters
because the optimal set of a convex problem is a slice
$\{x\in P:Nx=Np\}$ with a known rational matrix $N$ but an unknown, possibly
irrational, right-hand side $Np$.

\begin{lemma}[Hoffman bound with explicit constant]\label{lem:integer-hoffman}
Let $A\in\R^{k\times s}$, $T\in\R^{l\times s}$, $b\in\R^k$, and $t\in\R^l$,
and let $Q=\{x:Ax\le b\}$ and $Z=\{x\in Q:Tx=t\}$ with $Z\ne\emptyset$. Let
$C\ge1$ bound the absolute values of the entries of $A$ and $T$, and let
$\mu\in(0,1]$ be such that every linearly independent family of rows of $A$
and $T$ has a maximal square submatrix with determinant of absolute value
at least $\mu$. Then
\[
 \dist(x,Z)\le\frac{(sC)^{s-1}}{\mu}\,\norm{Tx-t}\qquad(x\in Q).
\]
For integer matrices $A$ and $T$ the determinant condition holds with
$\mu=1$.
\end{lemma}

The third tool converts an effective error modulus into a fixed selector.
It is a quantitative form of Tikhonov regularization, which is known to
select the minimum-norm optimizer in the limit; the schedule below makes
the rate explicit and charges only the bit length of the parameter.

\begin{lemma}[Fixed minimum-norm selector]\label{lem:selector}
Let $P\subseteq\R^n$ be nonempty, closed, and convex, and let
$\varphi:P\to\R$ be convex and continuous with nonempty optimal set $S$ and
optimal value $\varphi_*$. Let $p$ be the point of $S$ of minimum norm.
Suppose that $d\ge2$, $\Gamma\ge1$, and $R\ge1$ satisfy $\norm p\le R$ and
\[
 \dist(x,S)\le\Gamma(\varphi(x)-\varphi_*)^{1/d}\qquad
 (x\in P,\ \varphi(x)-\varphi_*\le1).
\]
For $\varepsilon\in(0,1]$ let
$\tau=\varepsilon^{2d-2}/(4\cdot8^{d-1}R^d\Gamma^d)$. Then
$\Phi_\tau=\varphi+\tau\norm{\cdot}^2$ has a unique minimizer $x_\tau$ on
$P$, with $\norm{x_\tau}\le\norm p$ and $\norm{x_\tau-p}\le\varepsilon/2$,
and every $x\in P$ with $\Phi_\tau(x)-\Phi_\tau(x_\tau)\le\tau\varepsilon^2/4$
satisfies $\norm{x-p}\le\varepsilon$. If every point of $S$ has norm at most
$R$, the factor $4$ in $\tau$ may be replaced by $2$.
\end{lemma}

The optimal set $S$ need not be bounded; the projection of $x_\tau$ onto
$S$ may lie far outside the ball of radius $R$, and the factor $4$ accounts
for it. Since $\log_2(1/\tau)=(2d-2)q+O(d)+d\log_2R+d\log_2\Gamma$, the
selector costs polynomially many value bits exactly when $\log\Gamma$ and
$\log R$ are polynomial. The remaining question in this section is when they
are.

\subsection{Globally convex sparse polynomials}\label{sec:points-global}

\begin{theorem}[Global point oracle]\label{thm:global-point}
Let $f\in\Q[x_1,\ldots,x_n]$ be an explicit sparse polynomial that is
convex on $\R^n$, and let $P=\{x\in\R^n:Ax\le b\}$ with rational $A$ and $b$.
Let $L$ be the encoding length of $(f,A,b)$ and $D=\max\{2,\deg f\}$. There
is a deterministic algorithm with the following properties. All running
times are bounded by polynomials, with absolute exponents, in the indicated
quantities, uniformly in $D$.
\begin{enumerate}[label=(\alph*)]
\item In time polynomial in $L,D$ it decides whether $P$ is empty and, if
not, whether $f$ is unbounded below on $P$. In the unbounded case it
returns rational $a\in P$ and $d$ with $Ad\le0$ and $f(a+td)=f(a)-t$ for all
$t\ge0$.
\item Otherwise $f$ attains its minimum $f_*$ on $P$, and the optimal set
$S$ has a unique point $p$ of minimum Euclidean norm. In time polynomial in
$L,D$ the algorithm computes rationals $R\ge1$ and $\Gamma\ge1$, with
$\log R+\log\Gamma$ bounded by a polynomial in $L,D$, such that
$\norm p\le R$ and
\[
 \dist(x,S)\le\Gamma\bigl(f(x)-f_*\bigr)^{1/D}\qquad
 (x\in P,\ f(x)-f_*\le1).
\]
\item For every integer $q\ge0$ it computes, in time polynomial in
$L,D,q$, a rational $x_q\in P$ with $\norm{x_q-p}\le2^{-q}$. For every further
integer $q_f\ge0$ it can return, in time polynomial in $L,D,q+q_f$, also a
rational $\ell$ with $\ell\le f_*\le f(x_q)\le\ell+2^{-q_f}$.
\end{enumerate}
\end{theorem}

The bounds are polynomial in the numerical degree $D$. They are polynomial
in the input length when exponents are written in unary, or more generally
when $D$ is polynomially bounded in $L$. The error bound holds on the whole
polyhedron $P$, which may be unbounded, and not only on a bounded part of it.
The selector $p$ depends only on the instance, so the outputs for
different $q$ approximate the same point. The decision in (a), attainment,
and a norm bound of polynomial bit length for some optimizer were proved by
\citet{SlotSteurerWiedmer2025} for inputs with unary degrees; we rederive
them with explicit constants. The effective modulus and the fixed selector
are the new outputs.

\paragraph{Proof outline.}
The complete proof is in Appendix~\ref{app:points-global}. It has four
steps.

\emph{Step 1: a computable quadratic minorant.} For every $a,x$, univariate
interpolation of $f$ along the line through $a$ and $x$ gives
\[
 f(x)\ge f(a)+\nabla f(a)^{\mathsf T}(x-a)+\frac{(x-a)^{\mathsf T}\nabla^2f(a)
 (x-a)}{\kappa_D},\qquad\kappa_D=(D+1)D^{D+2}.
\]
Averaging over $a$ in the unit cube yields
$f(x)\ge m_0+g^{\mathsf T}x+x^{\mathsf T}Mx/\kappa_D$, where
$M=\int\nabla^2f(a)\,da$ and the other coefficients are rational and
computed monomial by monomial. The kernel of $M$ is the common kernel of all
Hessians. With $\Pi$ the projection onto the range of $M$ and
$w=-(I-\Pi)\nabla f(0)$,
\[
 f(x)=f(\Pi x)-w^{\mathsf T}x\qquad(x\in\R^n),
\]
so $f$ is a function of $\Pi x$ plus a linear term in the remaining
directions.

\emph{Step 2: boundedness and a radius.} One linear system,
$w=A^{\mathsf T}\lambda+Mz$ with $\lambda\ge0$, decides boundedness. If it is
infeasible, Farkas' lemma gives a recession direction of $P$ along which $f$
decreases linearly. If it is feasible,
$f(x)\ge\beta_0-\nu\norm{\Pi x}+\mu_0\norm{\Pi x}^2$ on $P$ for explicit
rationals $\beta_0$, $\nu\ge0$, and $\mu_0>0$. On a sublevel set this bounds $\norm{\Pi
x}$, and together with convexity at the origin it bounds $w^{\mathsf T}x$ on
both sides; the constraints $Ax\le b$ alone give only one side. Points with
equal values of $Mx$ and $w^{\mathsf T}x$ have equal objective values, so
Lemma~\ref{lem:integer-hoffman} replaces every point of the sublevel set by
an equally good feasible point of norm less than an explicit $R$. This
proves attainment and $\norm p<R$ without assuming either.

\emph{Step 3: an effective modulus.} Let $x\in P$ have gap $E\le1$ and let
$h(u)=f(p+u)-f(p)-\nabla f(p)^{\mathsf T}u$, which is convex and nonnegative
on $\R^n$. On the segment from $p$ to $x$ we have $0\le h\le E$, and
interpolation extrapolates this bound along the whole line. For a point $c$
of the unit cube, convexity of $h$ at the midpoint of $2(c-p)$ and
$2t(x-p)$ bounds $h(c-p+t(x-p))$ for $0\le t\le E^{-1/D}$; differentiating at
$t=0$ gives
\[
 |\nabla f(c)^{\mathsf T}(x-p)|\le C_*E^{1/D}\qquad(c\in[0,1]^n),
\]
with an explicit $C_*$ that depends on $\norm p$ but not on $x$. The
polynomial $z\mapsto\nabla f(z)^{\mathsf T}(x-p)$ has coefficient vector
$N(x-p)$, where the rows of the rational matrix $N$ are read from the
monomials of $\nabla f$. A tensor interpolation inequality bounds these
coefficients by $\beta_{D-1}^nC_*E^{1/D}$, where $\beta_k=(k+1)2^kk^k$; the
interpolation grid appears only in the proof. At zero gap this shows $S=\{x\in P:Nx=Np\}$, and
Lemma~\ref{lem:integer-hoffman} gives $\Gamma$.

\emph{Step 4: selection.} Lemma~\ref{lem:selector} with $d=D$ gives a
parameter $\tau$ with polynomially many bits. The regularized minimizer has
norm at most $\norm p<R$, so it is also the minimizer over the bounded
polytope $P\cap[-R,R]^n$, to which Lemma~\ref{lem:convex-value} applies.

\begin{corollary}[Bounded polytopes]\label{cor:points-bounded}
Suppose in addition that $P$ is bounded and a rational $\varrho\ge1$ with
$P\subseteq[-\varrho,\varrho]^n$ is given. Then parts (b) and (c) of
Theorem~\ref{thm:global-point} hold with $R=\max\{1,n\varrho\}$ and the
constant $\Gamma$ of Step~3, without Steps~1 and~2, and the selector may use
$\tau=\varepsilon^{2D-2}/(2\cdot8^{D-1}R^D\Gamma^D)$. In particular, for
every fixed degree bound, the minimum-norm optimizer of a globally convex
polynomial on a bounded rational polytope can be approximated to $2^{-q}$ in
time polynomial in $L+q$.
\end{corollary}

\begin{remark}[Representation and degree]\label{rem:points-expansion}
The algorithms evaluate $f$ at rational points, also after affine changes
of coordinates, and never expand $f(Tx+t)$. Expansion can be exponentially
long in the dimension even when the degree is written in unary: for
$n=2r\ge4$, the polynomial
$f(y)=\sum_iy_i^{2r}$, convex on $\R^n$, has $n$ monomials, whereas
$f(Uy)$ for the rational orthogonal matrix
$U=I-\frac2n\mathbf1\mathbf1^{\mathsf T}$ has at least $2^r$ monomials
(Appendix~\ref{app:points-global}). The theorem is polynomial in the
numerical degree $D$. With exponents written in binary, $D$ can be
exponential in $L$, and we claim no bound polynomial in $\log D$. The
theorem also does not apply to polynomials given by arithmetic circuits: the
matrix $N$ of Step~3 is read from the monomial list.
\end{remark}

The exponent $1/D$ cannot be improved in general. A polynomial of odd
degree at least three is not convex on $\R^n$, because its restriction to a
suitable line has a second derivative of odd degree, which takes negative
values. For even $D$,
$f(t)=t^D$ on $P=\R$ has $\dist(t,S)=(f(t)-f_*)^{1/D}$.

\subsection{Cubics convex on a polytope}\label{sec:points-cubic}

A polynomial of degree at most three that is convex on $\R^n$ is quadratic
(Lemma~\ref{lem:models-low-degree}), and for such polynomials exact convex
quadratic programming already returns rational optimizers
\citep{KozlovTarasovKhachiyan1980}. Genuine cubics are convex only on part of
space, and convexity on the feasible polytope is the natural promise; for
cubics on boxes it is strongly NP-hard to recognize
\citep[Theorem~2.3]{AhmadiHall2020}. The next theorem shows that this weaker
promise still yields an effective modulus.

\begin{theorem}[Cubics convex on a polytope]\label{thm:cubic-point}
Let $P=\{x\in\R^n:Ax\le b\}$ be a nonempty bounded rational polytope and let
$f\in\Q[x_1,\ldots,x_n]$ be an explicit polynomial of degree at most three
that is convex on $P$. Let $L$ be the encoding length of $(f,A,b)$ and
$S=\argmin_Pf$. In time polynomial in $L$ one can compute a rational
$\Gamma_P\ge1$ with $\log\Gamma_P$ polynomial in $L$ such that
\[
 \dist(x,S)\le\Gamma_P\bigl(f(x)-f_*\bigr)^{1/4}\qquad
 (x\in P,\ f(x)-f_*\le1).
\]
For every integer $q\ge0$, in time polynomial in $L+q$ one can compute a
rational $x\in P$ with $\dist(x,S)\le2^{-q}$, and also a rational $x_q\in P$
with $\norm{x_q-p}\le2^{-q}$, where $p$ is the point of $S$ of minimum
Euclidean norm. No uniqueness of the optimizer, strong convexity, or
interiority of the optimizer is assumed, and $P$ may be lower dimensional.
\end{theorem}

Two facts specific to degree three drive the proof. The Hessian
$H(u)=\nabla^2\varphi(u)$ of a cubic $\varphi$ is affine in $u$, and the
third derivative is a constant symmetric tensor. The first fact lets an
interior reference point control every Hessian on the polytope. The second
links the curvature at that point to the curvature at an optimizer and a
feasible point. The recourse analysis reuses the following two lemmas with a
box in place of the polytope.

\begin{lemma}[Cubic transverse estimate]\label{lem:points-cubic-transverse}
Let $\varphi$ be a real polynomial of degree at most three on $\R^m$, let
$\Omega\subseteq\R^m$ be convex with $\nabla^2\varphi\succeq0$ on $\Omega$,
let $v$ minimize $\varphi$ over $\Omega$, and let $u\in\Omega$. Put $d=u-v$
and $E=\varphi(u)-\varphi(v)$. For every $c\in\R^m$ and all $\Lambda$,
$\rho_1$, $\rho_2$ with
$\Lambda\ge\max\{\norm{\nabla^2\varphi(u)},\norm{\nabla^2\varphi(v)}\}$,
$\rho_1\ge\norm{c-v}$, and $\rho_2\ge\norm d$,
\[
 \bigl(d^{\mathsf T}\nabla^2\varphi(c)\,d\bigr)^2\le12\Lambda(\rho_1+\rho_2)^2E .
\]
\end{lemma}

\begin{lemma}[Cubic optimal slice]\label{lem:points-cubic-slice}
Let $\varphi$ and $\Omega$ be as in Lemma~\ref{lem:points-cubic-transverse},
let $c\in\Omega$, and suppose that
$K=\ker\nabla^2\varphi(c)\subseteq\ker\nabla^2\varphi(u)$ for every
$u\in\Omega$. Let $g_c=\nabla\varphi(c)$ and let $\Pi_K$ be the orthogonal
projection onto $K^\perp$. If $v$ minimizes $\varphi$ over $\Omega$, then
\[
 \argmin_\Omega\varphi=\{w\in\Omega:\nabla^2\varphi(c)(w-v)=0,\
 g_c^{\mathsf T}(w-v)=0\}.
\]
If moreover $\Lambda\ge\sup_\Omega\norm{\nabla^2\varphi}$ and
$R_\Omega\ge\sup_{u\in\Omega}\norm{u-c}$ are finite, then
$|g_c^{\mathsf T}(u-v)|\le\varphi(u)-\varphi(v)+\Lambda
R_\Omega\norm{\Pi_K(u-v)}$ for every $u\in\Omega$.
\end{lemma}

The gradient row in Lemma~\ref{lem:points-cubic-slice} is necessary: along
the common Hessian kernel a cubic can still change linearly.

\paragraph{Proof outline.}
The proof is in Appendix~\ref{app:points-cubic}. Linear programming finds
the affine hull of $P$, a rational parametrization $x=x_0+Vu$, and a ball of
radius $r_0$ inside the reduced polytope $P'$, which lies in a ball of
radius $R_0$ about the same center. This reduction must come first:
convexity on a lower-dimensional polytope does not make the ambient Hessian
positive semidefinite. Reflection through the center and affinity of the
Hessian give $0\preceq H(u)\preceq(1+R_0/r_0)H_c$ on $P'$, where $H_c$ is the
Hessian at the center. Hence $\ker H_c$ is the common kernel, it is rational,
and individual boundary Hessians may only have larger kernels.
Lemma~\ref{lem:points-cubic-transverse} gives
$E\ge(d^{\mathsf T}H_cd)^2/(108\Lambda R_0^2)$, which bounds the displacement
transverse to $\ker H_c$ by a multiple of $E^{1/4}$ once the least positive
eigenvalue of $H_c$ is bounded below by its principal minors.
Lemma~\ref{lem:points-cubic-slice} describes $S$ as a slice with the rational
matrix $\binom{H_c}{g_c^{\mathsf T}}$, and Lemma~\ref{lem:integer-hoffman}
gives $\Gamma_P$. Lemmas~\ref{lem:selector} and~\ref{lem:convex-value} then
give both outputs. The regularizer is the original norm $\norm x^2$, not the
norm of the reduced coordinates $u$, which would select a different optimizer
in general.

The exponent $1/4$ suffices for polynomial dependence on $q$. The cubic
$t^3$ on $[0,1]$ has $\dist(t,S)=(f(t)-f_*)^{1/3}$, so no exponent larger than
$1/3$ is possible; whether $1/3$ is attainable with polynomial-bit constants
in all dimensions is open.

\subsection{Where the conversion fails}\label{sec:points-limits}

The positive results use global convexity or degree three. The constructions
below show that neither hypothesis can be dropped for quartics, and that
outputs beyond a set approximation can be hard even when set approximations are
easy. They use two standard arithmetic decision problems. \emph{Square Root
Sum} asks, for positive integers $a_1,\ldots,a_n$ and an integer $B$, whether
$\sum_i\sqrt{a_i}\le B$.
\emph{$\PosSLP$} asks whether the integer computed by a straight-line program
with constants $0,1$ and operations $+,-,\times$ is positive.
\citet{AllenderEtAl2009} show that Square Root Sum is decidable in polynomial
time with a $\PosSLP$ oracle and that both problems lie in the counting
hierarchy; by Theorem~\ref{thm:posslp-closure}(b), that oracle algorithm can be
replaced by a polynomial-time many-one reduction to one $\PosSLP$ instance.
Whether either problem lies in $\mathrm P$ is open. The statements below are
reductions and conditional consequences, not NP-hardness results. Their
structural bounds refer to the \emph{interaction graph} of a polynomial,
whose vertices are the variables, two variables being adjacent when some
monomial with a nonzero coefficient contains both.

\begin{proposition}[Exact active bounds for strongly convex cubics]\label{prop:points-active}
There is a polynomial-time algorithm that maps every Square Root Sum
instance either to its answer or to an explicit rational cubic $F$ on a
rational box $X$ with a distinguished coordinate $\theta\in[0,1]$ such that
\begin{enumerate}[label=(\roman*)]
\item $\nabla^2F\succeq\frac1{16}I$ on $X$, so $F$ has a unique minimizer
$z^*$;
\item $\theta^*=0$ if and only if $\sum_i\sqrt{a_i}\le B$, and otherwise
$0<\theta^*<1$;
\item the box widths and the absolute values of the coefficients of $F$ are
bounded by an absolute constant, the diagonal entries of $\nabla^2F$ are at
most $5$ on $X$, and the interaction graph of $F$ has treewidth at most two.
\end{enumerate}
A point within $2^{-q}$ of $z^*$ is computable in time polynomial in $L+q$.
\end{proposition}

Thus exact active-bound recognition for uniformly strongly convex cubics
with bounded data and treewidth two is at least as hard as Square Root Sum,
including
the equality case, although Theorem~\ref{thm:cubic-point} and even plain
strong convexity give points. The base of the construction is a complete
binary averaging tree whose cubic leaves $u_i^3/3-c_iu_i$ select normalized
square roots and whose squared residuals propagate their average to the
root; a sign test on the root then activates or deactivates the
distinguished coordinate.

\begin{proposition}[Selection versus some optimizer]\label{prop:points-selector}
There is a polynomial-time algorithm that maps every Square Root Sum
instance either to its answer or to an explicit rational quartic $F'$ on a
rational box, convex on that box, with the same treewidth and coefficient
bounds, such that its optimal set is a segment $\{z^*\}\times[0,1]$ if
$\sum_i\sqrt{a_i}\le B$ and the single point $(z^*,1)$ otherwise. The
point $(z^*,1)$ is optimal in both cases and can be approximated to $2^{-q}$
in time polynomial in $L+q$, whereas the last coordinate of the minimum-norm
optimizer is $0$ in the first case and $1$ in the second.
\end{proposition}

A set approximation is therefore easy for this family, while the
minimum-norm selector at accuracy $1/4$, or the dimension of the optimal set,
decides Square Root Sum. Requiring minimum norm strengthens the output
contract. The next
theorem shows that for quartics convex only on a box even a set
approximation at constant accuracy can carry the arithmetic difficulty.

\begin{theorem}[Constant-accuracy points for box-convex quartics]\label{thm:points-quartic-lower}
\leavevmode
\begin{enumerate}[label=(\alph*)]
\item There is a polynomial-time algorithm that maps every Square Root Sum
instance either to its answer or to an explicit rational quartic $F$ on
$[0,1]^N$, convex on $[0,1]^N$, with a unique minimizer $z^*$ whose
designated coordinate is $1$ if $\sum_i\sqrt{a_i}>B$ and $0$ if
$\sum_i\sqrt{a_i}<B$; equality does not occur in the produced instances. The
interaction graph of $F$ has treewidth at most two, every coefficient other
than the constant term has absolute value bounded by an absolute constant,
and the diagonal Hessian entries are at most $20$ on $[0,1]^N$.
\item There is a polynomial-time algorithm that maps every straight-line
program, with output $A$, to an explicit rational quartic $F$ on $[0,1]^N$,
convex on $[0,1]^N$, with a unique minimizer whose designated coordinate is
$1$ if $A>0$ and $0$ if $A\le0$. All coefficients of $F$ have absolute value
bounded by an absolute constant.
\item Consequently, if some deterministic algorithm returns, for every
explicit rational quartic convex on a rational box with a unique minimizer,
a point within Euclidean distance $1/4$ of the minimizer in time polynomial
in the input length, then Square Root Sum and $\PosSLP$ lie in $\mathrm P$; for Square
Root Sum the instances of (a) suffice. If the algorithm is always correct and
runs in expected polynomial time, both problems have Las Vegas algorithms
with expected polynomial running time.
\end{enumerate}
\end{theorem}

Both constructions place two independent strongly convex sign tests side by
side, one activated when the arithmetic quantity is positive and the other
when it is negative, and couple their activation amplitudes $\theta_\pm$ to a
new coordinate $y\in[0,1]$ by
$\frac1{192}[\theta_+^2(y-1)^2+\theta_-^2y^2]$. The coupling is not convex by
itself, but the curvature of the sign tests absorbs its negative part, so the
whole objective is convex on the box. The coupling vanishes at the right
endpoint and forces every minimizer to it, however small the active
amplitude is. For Square Root Sum, an equality
$\sum_i\sqrt{a_i}=B$ forces every radicand to be a perfect square, by a trace
argument in the multiquadratic field (Lemma~\ref{lem:points-trace}), so
equality is decided by integer square roots. For $\PosSLP$, gate values are
encoded as bounded pairs $(u,v)$ with $u/v$ equal to the integer value, the
gate equations are squared with geometric weights that make the objective
strongly convex on its box, and the output signal $(2u-v)/4$ has the sign of
$2A-1\ne0$. The final instances are convex on their boxes but not on
$\R^N$, and, before any rescaling, their curvature in the direction of $y$ at
the minimizer is $\frac1{96}((\theta_+^*)^2+(\theta_-^*)^2)$, which can be
extremely small. So
the theorem is consistent with Theorem~\ref{thm:global-point}, with
Theorem~\ref{thm:cubic-point}, with polynomial-time value enclosures from
Lemma~\ref{lem:convex-value}, and with algorithms whose running time depends
on a supplied growth constant. \citet{EY2010} proved $\PosSLP$-hardness of
strong approximation of equilibrium coordinates, and \citet{TarasovVyalyi2008}
reduced arithmetic circuit comparison to exact semidefinite feasibility;
Theorem~\ref{thm:points-quartic-lower} transfers the first phenomenon to
minimizers of polynomials convex on a box.

The last result of this subsection is unconditional, but it concerns the
value-to-point conversion rather than every algorithm.

\begin{proposition}[Value tolerance and isotropic regularization]\label{prop:points-regularization}
For $n\ge1$ let
\[
 F_n(x,y)=8\sum_{i=1}^nx_i^2-\frac{x_1}4-\sum_{i=2}^nx_{i-1}^2x_i+x_n^2(y-1)^2
 \qquad\text{on }[0,1]^{n+1}.
\]
\begin{enumerate}[label=(\alph*)]
\item $F_n$ is a quartic with $O(n)$ monomials and constant-size rational
coefficients, convex on $[0,1]^{n+1}$, with the unique minimizer
$z^*=(x^*,1)$, where $0<x_n^*\le2^{-2^{n+1}}$.
\item The feasible point $(x^*,0)$ has objective gap at most
$2^{-2^{n+2}}$ and distance one from $z^*$. Hence every error bound
$\dist(z,S)\le\Gamma(F_n(z)-F_{n,*})^\vartheta$ valid for gaps at most one has
$\log_2\Gamma\ge\vartheta\,2^{n+2}$.
\item For $\lambda>0$ let $z_\lambda$ minimize $F_n+\lambda\norm z^2$ on
$[0,1]^{n+1}$. If $\norm{z_\lambda-z^*}\le\frac12$, then
$\lambda\le2^{-2^{n+2}}$, so the denominator of a rational $\lambda$ has more
than $2^{n+2}$ bits.
\item Nevertheless $F_n(z_\lambda)-F_{n,*}\le(n+1)\lambda$ for every
$\lambda>0$, and a point within $2^{-q}$ of $z^*$ is computable in time
polynomial in $n+q$.
\end{enumerate}
\end{proposition}

With every exponent vector written in full
(Definition~\ref{def:models-polynomial}), the input has length
$L=\Theta(n^2)$. The quantities bounded below in (b) and (c),
$(\log_2\Gamma)/\vartheta$ and the bit length of $\lambda$, are exponential in
$n$, hence of the form $2^{\Omega(\sqrt L)}$: superpolynomial, but not
exponential, in $L$. Part (b) says that a value tolerance with fewer than
$2^{n+2}$ bits does not certify any point accuracy below one. Every error
modulus of this family has $(\log_2\Gamma)/\vartheta\ge2^{n+2}$, so the value
precision $\vartheta^{-1}(q+\log_2\Gamma)$ needed to convert values into points
through it is exponential in $n$ whatever the exponent; for exponents bounded
below by an inverse polynomial, the constant itself has a number of bits
exponential in $n$. Theorem~\ref{thm:cubic-point} therefore does not extend to
quartics convex on a box, and the obstruction is in the modulus itself, not in
its proof. Part (c) shows that plain Tikhonov selection, which works in
Lemma~\ref{lem:selector} because $\Gamma$ is short, needs a parameter whose
bit length is exponential in $n$ here. Part (d) shows that this family is
nonetheless easy for both values and points: setting $y=1$ leaves a
uniformly strongly convex cubic. The difficulty is created by the method,
not by the instance.

\subsection{Representations and certificates}\label{sec:points-representation}

A procedure that returns rational approximations at every precision
describes an optimizer compactly, but it is not an exact representation. The
next two results show that the usual exact formats can have length
exponential in the dimension, hence superpolynomial in the input length,
even for strongly convex, uniformly conditioned instances whose interaction
graph is a path (or has treewidth two, for the multiplier rectangles), where
points are cheap.

\begin{proposition}[Expanded algebraic output on a conditioned path]\label{prop:points-algebraic-output}
For $n\ge2$ let
$\Psi_n(x)=\sum_{i=1}^{n-1}(x_i^2-x_{i+1}-2)^2+(x_n-1)^2$ on $[1,2]^n$.
\begin{enumerate}[label=(\alph*)]
\item $\Psi_n$ has the unique minimizer $a$ with $a_n=1$ and
$a_i=\sqrt{a_{i+1}+2}$. On $[1,2]^n$ it satisfies
$\Psi_n(x)\ge\frac14\norm{x-a}^2$ and $\partial_{ii}\Psi_n\le38$. Its
interaction graph is a path, and it has $O(n)$ monomials with constant-size
coefficients.
\item The minimal polynomial of $a_1$ over $\Q$ has degree $2^{n-1}$ and
exactly $2^{n-2}+1$ nonzero coefficients.
\item In time $O(n)$ one can compute a box $B=\prod_i[\ell_i,u_i]\subseteq
[1,2]^n$ with dyadic endpoints of $O(1)$ bits that contains $a$ and on which
$\nabla^2\Psi_n\succeq\frac1{32}I$.
\end{enumerate}
\end{proposition}

On the box $B$, $\Psi_n$ is a strongly convex quartic with a uniformly
conditioned path structure, and Lemma~\ref{lem:convex-value} with strong
convexity gives points within $2^{-q}$ of $a$ in time polynomial in $n+q$.
The recurrence $a_i=\sqrt{a_{i+1}+2}$ describes the whole optimizer with
$n-1$ additions and $n-1$ square-root gates. Yet a dense defining polynomial
of $a_1$ has more than $2^{n-1}$ coefficients, and the minimal polynomial
written as a sparse list has $2^{n-2}+1$ terms, while the input, with every
exponent vector written in full, has length $\Theta(n^2)$. The support bound
concerns the minimal polynomial; it does not exclude a sparser
non-minimal multiple. Theorem~\ref{thm:singleton-field} and the related
results on exact optimizers treat other formats.

\begin{proposition}[Rectangular certificates]\label{prop:points-rectangle}
For $n\ge1$ let $X=[0,\frac12]^{n+1}$, let
\[
 \Xi_n(x)=\Bigl(x_0-\frac14\Bigr)^2+\sum_{i=1}^n\Bigl(x_i-\frac{x_{i-1}^2}4\Bigr)^2,
\]
and let $m=4\cdot2^n-2$.
\begin{enumerate}[label=(\alph*)]
\item $\Xi_n$ has the unique minimizer $a$ with $a_i=2^{-(4\cdot2^i-2)}$,
which lies in the interior of $X$. On $X$, $\nabla^2\Xi_n\succeq\frac58I$,
$\Xi_n(x)\ge\frac9{16}\norm{x-a}^2$, all diagonal Hessian entries are at most
$\frac{35}{16}$, and the interaction graph is a path. Every rectangle with
rational endpoints that contains $a$ and lies in the open box
$(0,\frac12)^{n+1}$ has an endpoint with denominator at least $2^m$.
\item Let
$\widetilde\Xi_n(x,y)=\Xi_n(x)+y^2+\frac y4\bigl(x_n-\frac18x_{n-1}^2\bigr)$
on $X\times[0,\frac12]$. It has the unique minimizer $(a,0)$, satisfies
$\nabla^2\widetilde\Xi_n\succeq\frac{11}{32}I$ and
$\widetilde\Xi_n\ge\frac{135}{256}(\norm{x-a}^2+y^2)$, has diagonal Hessian
entries at most $\frac{35}{16}$ and treewidth two, and its active lower bound
$y=0$ has $\partial_y\widetilde\Xi_n(a,0)=a_n/8>0$. Every rectangle
$B\subseteq X$ with rational endpoints that contains $a$ and satisfies
$\partial_y\widetilde\Xi_n(x,0)\ge0$ for all $x\in B$ has an endpoint with
denominator at least $2^m$.
\item Both optimizers have exact descriptions of size $O(n)$: the
recurrence $a_0=\frac14$, $a_i=\frac14a_{i-1}^2$ as an arithmetic circuit,
with optimality certified by a contraction argument on the closed box $X$
and, in (b), by the displayed lower bound.
\end{enumerate}
\end{proposition}

The two rectangle formats are natural for interval methods: a box strictly
inside the feasible region that encloses the optimizer, and a box on which
the sign of an active multiplier is certified uniformly. Under uniform
conditioning, fixed degree, and treewidth at most two, both can require
endpoints whose bit length is exponential in $n$, while correlated implicit
certificates, which keep the residual equations, stay short. These are lower
bounds for specified formats, not for exact optimization in general; in
particular, boxes that touch the boundary
of $X$ and certificates of interval-Newton type that use the residual
equations are not excluded.

\subsection{Summary}\label{sec:points-summary}

Table~\ref{tab:points-summary} collects the results by output contract.
Values are available in all three settings under a convexity promise. The
route from values to set approximations goes through an effective error
modulus. It exists under global convexity and for cubics convex on a
polytope, but not in general for quartics convex on a box, where even
constant-accuracy points are as hard as Square Root Sum or $\PosSLP$. A
fixed selector needs, in addition, a selection rule; the minimum-norm rule
costs nothing extra when the modulus is effective
(Lemma~\ref{lem:selector}), but it can be hard where a set approximation is
easy (Proposition~\ref{prop:points-selector}). Exact labels and exact
representations are different problems: they can be hard, or long, for
strongly convex cubics and quartics whose points are easy. Exact comparisons
under global strong convexity are treated in Theorem~\ref{thm:exact-upper}.

\begin{table}[ht]
\centering
\small
\begin{tabularx}{\linewidth}{@{}>{\raggedright\arraybackslash}p{0.2\linewidth}
>{\raggedright\arraybackslash}X>{\raggedright\arraybackslash}X
>{\raggedright\arraybackslash}X@{}}
\toprule
Output & Convex on $\R^n$, sparse, degree $D$ & Cubic, convex on a polytope &
Quartic, convex on a box\\
\midrule
Value enclosure & polynomial (Theorem~\ref{thm:global-point}) & polynomial
(Lemma~\ref{lem:convex-value}) & polynomial (Lemma~\ref{lem:convex-value})\\
Effective modulus & exponent $1/D$, polynomial bits
(Theorem~\ref{thm:global-point}) & exponent $1/4$, polynomial bits
(Theorem~\ref{thm:cubic-point}) & none for some families: every modulus has
$(\log\Gamma)/\vartheta\ge2^{n+2}$
(Proposition~\ref{prop:points-regularization})\\
Set approximation & polynomial in $L,D,q$ & polynomial in $L+q$ & decides
Square Root Sum or $\PosSLP$ at accuracy $1/4$
(Theorem~\ref{thm:points-quartic-lower})\\
Minimum-norm selector & polynomial in $L,D,q$ & polynomial in $L+q$ &
decides Square Root Sum where some optimizer is easy
(Proposition~\ref{prop:points-selector})\\
Exact active label & an exact comparison, not decided by points
(Theorem~\ref{thm:exact-upper} treats the strongly convex unconstrained
case) & decides Square Root Sum, even with uniform strong convexity
(Proposition~\ref{prop:points-active}) & as for cubics\\
Expanded output & not implied by points & not implied by points &
minimal polynomials and rectangle endpoints exponential in $n$ for strongly
convex instances of treewidth at most two
(Propositions~\ref{prop:points-algebraic-output}
and~\ref{prop:points-rectangle})\\
\bottomrule
\end{tabularx}
\caption{Output contracts for convex polynomial minimization in the bit
model. ``Polynomial'' means time polynomial in $L,D,q$ in the global convex
column and in $L,q$ in the fixed-degree columns, under the stated promises.
Hardness entries are reductions, not unconditional lower bounds,
except for the bit counts of moduli and formats.}
\label{tab:points-summary}
\end{table}
